\section{Setting and statements}\label{sec:statements}

This section fixes the objects and states the results. Sections~\ref{sec:base}--\ref{sec:examples} use the
definitions made here without repeating them.

\subsection{Conventions}\label{ssec:conventions}
We write $N\cong\Z^4$ and $M=\operatorname{Hom}(N,\Z)$, with pairing $\langle\ ,\ \rangle$. The polytope
$\Delta\subset N_\R$ is reflexive, with polar $\Delta^\circ=\{m\in M_\R:\langle m,n\rangle\ge-1\text{ for all }n\in\Delta\}$.
For a face $G$ of $\Delta$ the dual face is $G^\circ=\{m\in\Delta^\circ:\langle m,n\rangle=-1\text{ for all }n\in G\}$; for a
two-face $G$ it is an edge, of lattice length $\ell(G^\circ)$. We write $\mathbb P_\Delta$ for the toric variety of the fan
over the faces of $\Delta$. Its anticanonical hypersurfaces have Newton polytope $\Delta^\circ$, and such a hypersurface is
\emph{$\Delta^\circ$-regular} if it meets every torus orbit transversally \cite[\S3]{Batyrev}. The Batyrev mirror family
consists of the anticanonical hypersurfaces with Newton polytope $\Delta$ in a crepant resolution $X_{\mathcal T}$ of
$\mathbb P_{\Delta^\circ}$ (Section~\ref{ssec:mirrorfamily}). In the notation of Gross \cite{Gross2005} our pair
$(\Delta^\circ,\Delta)$ is $(\Delta,\nabla)$. We call the polytope whose lattice points index the monomials of a
hypersurface its \emph{Newton side}, and the polytope whose subdivision gives the fan of the ambient toric variety its
\emph{fan side}; for $X$ these are $\Delta^\circ$ and $\Delta$, for the mirror family $\Delta$ and $\Delta^\circ$.
We write $dP_k$ for a del Pezzo surface of degree $k$.

For a finite set of lattice points $n$ with heights $H(n)$ the tropical polynomial is
$m\mapsto\min_n\bigl(H(n)+\langle m,n\rangle\bigr)$, $m\in M_\R$, and its tropical hypersurface is the locus where the minimum is
attained at least twice. This convention is compatible with the signs of \cite{Gross2005}, and it is used in
Sections~\ref{sec:base}--\ref{sec:legendre} and~\ref{sec:resolution}; Sections~\ref{sec:regime}--\ref{sec:periods} use $p=-m$ and maxima of valuations, as stated
there.

\subsection{Admissible polytopes and profiles}\label{ssec:admissible}
A \emph{$dP_7$ pentagon} is a lattice pentagon with exactly one interior lattice point, its \emph{centre} $o$, and no
lattice points on its boundary besides its vertices; a \emph{$dP_6$ hexagon} is a lattice hexagon with the same two
properties. Up to $GL_2(\Z)$ and translation the hexagon has vertices $w_0,\dots,w_5=(1,0),(1,1),(0,1),(-1,0),(-1,-1),(0,-1)$ and centre
$o=(0,0)$, and $w_j+w_{j+2}=o+w_{j+1}$ for all $j$ (indices modulo $6$).

\begin{definition}[admissible polytopes]\label{def:admissible}
The reflexive polytope $\Delta$ is \emph{admissible} if every edge of $\Delta$ is primitive, every two-face is a unimodular
triangle, a unit square, a $dP_7$ pentagon or a $dP_6$ hexagon, and $\ell(G^\circ)=1$ for every two-face $G$ that is not a
triangle. This is \cite[Def.~8.1]{Paper4}.
\end{definition}

The lattice points of an admissible $\Delta$ that lie on faces of dimension at most two are the vertices and the
centres of the pentagons and hexagons; we denote this set by $R$. A singular two-face $G$ of $\Delta$ (a square, a
pentagon or a hexagon) gives one singular point of $X$, on the torus orbit of the cone over $G$: a node, a cone over
$dP_7$, or a cone over $dP_6$ respectively (Section~\ref{ssec:partialres}).

\begin{definition}[profiles and node parallelograms]\label{def:profile}
A \emph{profile} $\sigma$ assigns to every singular two-face of $\Delta$ a subdivision into unit parallelograms and
unimodular triangles, as follows. A unit square is not subdivided. A pentagon receives its unique subdivision into two
unit parallelograms and one triangle, all containing the centre (Section~\ref{ssec:faces}). A hexagon receives either
\begin{itemize}
\item an \emph{$A_1$ subdivision}: two parallelograms $\{o,w_j,w_{j+1},w_{j+2}\}$, $\{o,w_{j+3},w_{j+4},w_{j+5}\}$ and the
two triangles $\{o,w_{j+2},w_{j+3}\}$, $\{o,w_{j+5},w_j\}$, for one $j\in\{0,1,2\}$; or
\item an \emph{$A_2$ subdivision}: the three parallelograms $\{o,w_j,w_{j+1},w_{j+2}\}$ with $j\equiv j_0\pmod 2$, for one
$j_0\in\{0,1\}$.
\end{itemize}
The parallelograms of these subdivisions are the \emph{node parallelograms}; $Q$ denotes their set.
We say that a hexagon is on the \emph{$A_1$ branch} or the \emph{$A_2$ branch} of $\sigma$ according to its
subdivision. The names refer to the root sublattices generated by differences of exceptional curve classes
on the $dP_6$ surface over its centre (Remark~\ref{rem:tie}).
\end{definition}

The vertices of a node parallelogram are labelled $a,b,c,d$ in cyclic order with $a+c=b+d$, and one of its diagonals,
$\{a,c\}$, is chosen; in a pentagon or hexagon we choose the diagonal through the centre. The \emph{circuit} of $q\in Q$
and the \emph{relation lattice} are
\[
\circv_q:=e_b+e_d-e_a-e_c\in\Z^R,\qquad K:=\Bigl\{\lambda\in\Z^Q:\ \sum_q\lambda_q\,\circv_q=0\Bigr\}.
\]
$K$ is saturated in $\Z^Q$. Changing the chosen diagonal at $q$ changes the signs of $\circv_q$, of the coordinate
$\lambda_q$ and of every node coefficient $\rho_q$ below. The results hold for every choice of diagonals, with $K$, the node
coefficients and the orientations of the vanishing spheres (Definition~\ref{def:orient}) all formed from that choice;
changing the diagonal at $q$ reverses the orientation of $\delta_q$, so $\rho_q\,\delta_q$ does not depend on it
(Lemma~\ref{lem:diag}).

\begin{hypothesis}{(Proj)}{projectivity}\label{hyp:P}
There is a projective toric variety $P'_\sigma\to\mathbb P_\Delta$ whose fan refines the fan over the faces of $\Delta$,
has its rays generated by lattice points of $\partial\Delta$, and induces the profile $\sigma$ on every two-face of $\Delta$.
\end{hypothesis}

Under~\ref{hyp:P} the cones of the fan of $P'_\sigma$ cut $\partial\Delta$ into lattice polytopes; this subdivision is
denoted $\Sigma'$. Let $X\subset\mathbb P_\Delta$ be a $\Delta^\circ$-regular anticanonical hypersurface and $X'_\sigma$ its
proper transform in $P'_\sigma$. Then $X'_\sigma\to X$ is crepant, $X'_\sigma$ has exactly one singular point for each
node parallelogram, an ordinary double point, and we write $q\in Q$ also for this node. A choice of diagonal at every
node parallelogram defines a small resolution $\widehat X\to X'_\sigma$, with exceptional curves $C_q\cong\mathbb P^1$;
for suitable choices, including the diagonals through the centres at the pentagons and hexagons, $\widehat X$ is
projective, induced by a regular refinement of $\Sigma'$ (Proposition~\ref{prop:partialres}). For every small resolution the
integral relations among the classes $[C_q]$ in $H_2(\widehat X;\Z)$ are exactly $K$, with $K$ formed from the diagonals
of $\widehat X$; moreover $H_2(\widehat X;\Z)$ is torsion-free (Propositions~\ref{prop:partialres} and~\ref{prop:integral}). In particular, for a projective $\widehat X$ the number
$|Q|-\operatorname{rank}K$ is the relative Picard number of $\widehat X\to X'_\sigma$; it does not depend on the choice.

Hypothesis~\ref{hyp:P} is a linear feasibility problem in the heights of the lattice points of $\partial\Delta$. When
every singular two-face is a square it holds with $P'_\sigma=\mathbb P_\Delta$. It is a genuine restriction in general;
see Section~\ref{ssec:scope}.

\subsection{Gross's base and relative tropical \texorpdfstring{$2$}{2}-cycles}\label{ssec:tropical}
Let $\mathcal T$ be a unimodular triangulation of $\partial\Delta^\circ$ refining its faces, and $\check h$ an integral
height on the fan over $\mathcal T$ such that $\check h$ and $\check h-\check\varphi$ are strictly convex on that fan,
where $\check\varphi$ equals $1$ on $\partial\Delta^\circ$ and is linear on the cones over the faces of $\Delta^\circ$.
Such a pair exists for every admissible $\Delta$ (Theorem~\ref{thm:U}); we fix one, normalised as in
Section~\ref{ssec:gross}. The \emph{base} $B=B'_\sigma$ is the dual intersection complex, in the fan picture, of
Gross's toric degeneration of $X'_\sigma$ \cite{Gross2005}: the boundary of the polytope $\nabla^{\check h}$ with the
polyhedral decomposition $\mathscr P$ determined by $\Sigma'$ (Section~\ref{ssec:gross}). It is an integral affine
three-sphere with singularities along a trivalent-and-four-valent graph $\Gamma\subset B$, the \emph{discriminant}, which
has edges called \emph{arcs}. Each arc is a union of \emph{legs}, the segments joining the barycentres of an edge and
a two-cell of $\mathscr P$. Each leg $E$ carries a primitive \emph{vanishing vector} $d_E=n'-n\in N$, where
$\{n,n'\}$ is the lower support of the edge, namely the cell of $\Sigma'$ underlying it in Gross's construction
(Section~\ref{ssec:gross}); the
monodromy around $E$ is the transvection $x\mapsto x+\langle\alpha_E,x\rangle d_E$ for a covector $\alpha_E$. The legs
of an arc $A$ have the same lower support, so their vanishing vectors agree up to sign; we write $d_A$ for one of them
(Section~\ref{ssec:coef}). The
four-valent vertices of $\Gamma$ are in bijection with $Q$: the vertex $p_q$ lies on the two-cell of $\mathscr P$ over the
node parallelogram of $q$, and after an isotopy supported in that two-cell it is a negative node in the sense of
Casta\~no-Bernard and Matessi \cite[Def.~6.3]{CBM} (Proposition~\ref{prop:base}). We write $\Lambda$ for the local system of integral tangent vectors on $B\setminus\Gamma$ and
$\iota\colon B\setminus\Gamma\to B$ for the inclusion.

\begin{definition}[relative tropical $2$-cycles]\label{def:relcycle}
Let $\mathfrak K$ be a triangulation of $B$ in which $\Gamma$ is a full subcomplex; the barycentric subdivision of
$\mathscr P$ is one (Lemma~\ref{lem:full}). For a simplex $s$ of $\mathfrak K$ let $F(s)$ be the group of flat sections of
$\Lambda$ over $V\setminus\Gamma$ for a small neighbourhood $V$ of the closed simplex $\bar s$; faces restrict. Let
$C_*(\mathfrak K;F)$ be the simplicial chain complex with these coefficients, and $C_*(\Gamma;\mathcal D)\subset
C_*(\mathfrak K;F)$ the subcomplex of chains on $\Gamma$ whose coefficient on every edge $E$ of $\mathfrak K$ in $\Gamma$ lies in $\Z d_E$,
where $d_E:=\pm d_A$ for the arc $A$ whose closure contains $E$ (the sign does not affect $\Z d_E$). A
\emph{relative tropical $2$-cycle} is a chain $z\in C_2(\mathfrak K;F)$ with $\partial z\in C_1(\Gamma;\mathcal D)$. We
write $Z(\mathfrak K)$ for the group of relative tropical $2$-cycles, and call
$\partial z=\sum_Ec_E\,d_E\,[E]$ the \emph{boundary network} of $z$.
\end{definition}

In words: a relative tropical $2$-cycle is an integral $2$-chain on $B$ with coefficients in the sheaf $\iota_*\Lambda$
whose boundary lies on the discriminant, and runs along each leg in the direction of its vanishing vector. The
coefficient of a triangle is an integral tangent vector of $B$ that is invariant under the local monodromy wherever the
triangle meets $\Gamma$. The tropical $2$-cycles of Casta\~no-Bernard and Matessi \cite[Def.~7.2]{CBM}, triangulated, are
relative tropical $2$-cycles (Section~\ref{ssec:cbm}); we call the cycles of their definition \emph{strict}.
Here and below \cite[Def.~7.2]{CBM} is applied to $B'_\sigma$ with the two readings of Section~\ref{ssec:cbm}: clause~(i)
is read as conjugacy in $GL_2(\Z)$, and a negative vertex means a vertex of negative type
(Definition~\ref{def:types}); Theorem~\ref{thm:res} below holds also under the stricter reading of~(i) as conjugacy in
$SL_2(\Z)$, whose class of cycles is contained in this one, while the content of Conjecture~\ref{conj:strict} depends on
the reading.
The definition of Casta\~no-Bernard and Matessi allows a topological embedding $j$; throughout, strict cycles are taken with $j$ piecewise linear,
which the triangulation requires.
The relative ones are allowed non-primitive fields, arbitrary
branching along edges, crossings of $\Gamma$ with any invariant field, and boundary networks that pass a node in both
directions or with multiplicity.

The boundary network $\partial z$ is a $1$-cycle of $C_*(\Gamma;\mathcal D)$; such cycles are the \emph{balanced
networks}, and $\mathrm{Bal}(\Gamma)$ denotes their group. For a balanced network $\beta=\sum_Ec_E\,d_E\,[E]\in\mathrm{Bal}(\Gamma)$ let
$\tilde\beta=\sum_Ec_E(e_{n'}-e_n)[E]\in C_1(\Gamma;\Z^R)$, where $d_E=n'-n$, be its \emph{canonical lift}. The boundary of
$\tilde\beta$ vanishes away from the nodes, and at $p_q$ it is a multiple of the circuit (Lemma~\ref{lem:balanced}).

\begin{definition}[node coefficients]\label{def:nodecoef}
For $\beta\in\mathrm{Bal}(\Gamma)$ and $q\in Q$ the \emph{node coefficient} $\rho_q(\beta)\in\Z$ is defined by
$\partial\tilde\beta(p_q)=-\rho_q(\beta)\,\circv_q$. For $z\in Z(\mathfrak K)$ put $\rho(z):=\rho(\partial z)\in\Z^Q$.
\end{definition}

Equivalently, the four legs at $p_q$ have lower supports $\{a,b\}$, $\{b,c\}$, $\{c,d\}$ and $\{d,a\}$; if the
coefficients of a balanced network on the legs over $\{a,b\}$ and $\{b,c\}$ are oriented away from $p_q$ and written
$c_{ab}(b-a)$ and $c_{bc}(c-b)$, then $\rho_q=c_{ab}-c_{bc}$
(Section~\ref{ssec:cycles-bdry}). For a tropical $2$-cycle $S$ in the sense of \cite{CBM}, $\rho(S)\in\{0,\pm1\}^Q$
records the jump of the boundary field of $S$ across each node that $\partial S$ meets.

\subsection{The mirror family and its vanishing spheres}\label{ssec:mirrorfamily}
Let $X_{\mathcal T}$ be the smooth projective toric fourfold of the fan over $\mathcal T$, and put $\mathcal A=\Delta\cap N$.
The members of the mirror family are
\[
Y_{t,c}=\Bigl\{\sum_{n\in\mathcal A}c_n\,t^{H(n)}\,x^n=0\Bigr\}\subset X_{\mathcal T},\qquad 0<t<1,
\]
where $x^n$ is the Cox monomial of $n$, $c=(c_n)$ is a coefficient vector, and the height is $H=k_0\varphi+h_C$: here
$\varphi$ equals $1$ on $\partial\Delta$ and is linear on the cones over the faces of $\Delta$, $h_C$ is a strictly
convex height with integral slopes inducing the \emph{pulling refinement} $\mathcal F_C$ of $\Sigma'$
(Section~\ref{ssec:mirrorsub}), and $k_0$ is an integer bounded below in terms of the lattice data and $h_C$ (Lemma~\ref{lem:gaps}). As
$t\to0$ the members approach the maximally degenerate limit.

For a node parallelogram $P=abcd$ the \emph{node binomial ratio} of $c$ is $\mu_P=c_ac_c/(c_bc_d)$; the face polynomial
$c_ax^a+c_bx^b+c_cx^c+c_dx^d$ is a product of two binomials exactly when $\mu_P=1$. Let $\mathcal Z$ be the set of real
coefficient vectors with $c_0>0>c_n$ for $n\neq0$ and $\mu_P=1$ for every $P$. Every real $c$ with this sign pattern is
uniquely $c=c^{(0)}e^{\xi}$ with $c^{(0)}\in\mathcal Z$ and $\xi\in\operatorname{span}_\R\{r_P\}$, where
$r_P=e_a+e_c-e_b-e_d$, so that $r_P=-\circv_q$ under $\Z^R\subseteq\Z^{\Delta\cap N}$; then $\log\mu_P(c)=r_P\cdot\xi$. For $t$ small and $c$ in the window of Theorem~\ref{thm:main} below,
in which every $|\log\mu_P|$ lies between two fixed small positive bounds, the defining section of $Y_{t,c}$,
trivialised by dividing by its term $c_at^{H(a)}x^a$, has, for each $P$, a unique nondegenerate critical point $q_P$ in a polydisc $\mathbb D_P$ about a point of the torus orbit
associated with the edge dual to $P$ (Section~\ref{ssec:morse}). Its critical value for this trivialisation is small and nonzero, and $\delta_q\subset Y_{t,c}$
denotes its vanishing $3$-sphere, oriented by the convention of Definition~\ref{def:orient}.

\subsection{Results}\label{ssec:results}

\begin{mainthm}[relations among the vanishing spheres]\label{thm:main}
Let $\Delta$ be admissible, with a profile $\sigma$ satisfying~\ref{hyp:P}. There are $\eta_1>0$ and $C\ge1$, depending
only on the vectors $r_P$, $P\in Q$, and the auxiliary direction $\bar\xi$ fixed in Lemma~\ref{lem:eta}, with the
following property. For every $\eta_0\in(0,\eta_1]$
and every compact set $\mathscr C_0\subset\mathcal Z$ there is $t_0>0$ such that for $0<t<t_0$ and every real coefficient
vector $c$ with $c_0>0>c_n$ $(n\neq0)$, component $c^{(0)}\in\mathscr C_0$, and
$\eta_0/C\le|\log\mu_P(c)|\le\eta_0$ for every $P$, the member $Y_t=Y_{t,c}$ is smooth and:
\begin{enumerate}
\item \emph{(Relations.)} For $k\in\Z^Q$, the class $\sum_qk_q[\delta_q]$ is a torsion element of $H_3(Y_t;\Z)$ if and
only if $k\in K$, and it is then zero. So the relation lattice
$\Rel(\delta)=\{k:\sum_qk_q[\delta_q]=0\}$ equals $K$, and the classes $[\delta_q]$ span a free subgroup of
$H_3(Y_t;\Z)$, isomorphic to $\Z^Q/K$, of rank $|Q|-\operatorname{rank}K$, which meets the torsion subgroup only in $0$.
\item \emph{(Realisation.)} Every $k\in K$ is the vector of node coefficients $\rho(z)$ of a relative tropical
$2$-cycle $z$ on $B'_\sigma$, and $z$ lifts to an integral $4$-chain $\check\Xi(z)$ on $Y_t$ with
$\partial\check\Xi(z)=\sum_qk_q[\delta_q]$.
\item \emph{(The nodal locus.)} In an open neighbourhood $U$ of $c$ in the space of complex coefficient vectors
(Section~\ref{ssec:U}), the members with a node at every point $q_P$, $P\in Q$, form a nonempty complex submanifold of
codimension $|Q|-\operatorname{rank}K$, the relative Picard number of a projective small resolution
$\widehat X\to X'_\sigma$.
\end{enumerate}
\end{mainthm}

Morrison proposed that mirror symmetry exchanges the two sides of a conifold transition, with the number of conditions
imposed by the nodes on one side equal to the relative Picard number on the other \cite[\S3.1]{Mor99}. His setting
is a conifold transition, so $X'_\sigma$ is assumed smoothable there. Item~(3) is this count in the space of
coefficient vectors of the Batyrev family, for smoothable and non-smoothable $X'_\sigma$ alike. The factoring locus
of Section~\ref{ssec:periodsremarks}, cut out in $U$ by the equations $\mu_P=1$, that is $\eta_P=r_P\cdot\xi=0$ with
$\xi\in\C^{\mathcal A}$ as in Definition~\ref{def:window}, has the same codimension, since the $r_P$ span a space of
dimension $|Q|-\operatorname{rank}K$; item~(3) concerns the nodal locus itself, on which
all $|Q|$ nodes are present. Item~(1), with Proposition~\ref{prop:integral}, identifies the relations on the two sides
integrally: the exceptional curves of $\widehat X$ and the vanishing spheres of $Y_t$ have the same relation lattice $K$. For a projective conifold transition
the relations among the vanishing spheres and among the exceptional curves are mutual annihilators
\cite[Thm.~1.14]{LLW18}; item~(1) is the statement that crosses the mirror.

\begin{maincor}[smoothings and mirror relations]\label{cor:smooth}
Under the hypotheses of Theorem~\ref{thm:main} the following are equivalent:
\begin{enumerate}
\item[(a)] $X'_\sigma$ admits a smoothing;
\item[(b)] the exceptional curves satisfy a relation $\sum_q\lambda_q[C_q]=0$ in $H_2(\widehat X;\Q)$ with every
$\lambda_q\neq0$;
\item[(c)] the vanishing spheres satisfy a relation $\sum_qk_q[\delta_q]=0$ in $H_3(Y_t;\Z)$ with every $k_q\neq0$.
\end{enumerate}
The spheres $\delta_q$ are smoothly isotopic to pairwise disjoint Lagrangian spheres $L_q$ for the restriction of a
toric K\"ahler form. If (a)--(c) hold, one of the $2^{|Q|}$ conifold transitions of $Y_t$ in the $L_q$ carries a
symplectic structure in which the exceptional curves are symplectic.
\end{maincor}

Here a \emph{smoothing} is a flat deformation, over a germ of a complex space, whose general fibre is smooth. Batyrev
and Kreuzer, for polytopes whose singular two-faces are squares, specialised the mirror family to the locus where the
face polynomials of the squares factor, expected the specialised members to have a node in each of the local toric charts they single out, and observed that a projective small
resolution of such a nodal mirror requires a relation among the vanishing spheres with all coefficients nonzero;
they did not determine whether such a relation exists \cite[\S3]{BK10}. For dual edges of lattice length one, the
case considered here, their local equation near the expected node omits the monomials of the cells next to the
square, which are linear in the normal coordinates. The nodal members of Theorem~\ref{thm:main}(3) approach the
factoring locus as $t\to0$. Our estimates do not place them on it, and in a local model at a square node the critical
value on the factoring locus is nonzero for general coefficients (Section~\ref{ssec:periodsremarks}). For them Corollary~\ref{cor:smooth} settles the question: a
relation with all coefficients nonzero exists exactly when the original side is smoothable. For hexagon-free
$\Delta$ satisfying~\ref{hyp:P}, the same relation criterion decides the smoothability of $X$ itself (Corollary~\ref{cor:mixed}). The equivalence of~(a)
and~(b) is Namikawa's form \cite[Thm.~2.5]{Nam02} of the criterion of Friedman \cite[Rem.~(4.5), Cor.~(4.7)]{Friedman}
(compare \cite[Thms.~1.2 and~1.3]{RT09}); the symplectic statement is \cite[Thm.~2.9]{STY}.

Theorem~\ref{thm:main} is assembled from two results, which are the first two steps of the proof; the third step is
a period computation.

\begin{mainthm}[node coefficients on the base]\label{thm:real}
Let $\Delta$ be admissible with a profile satisfying~\ref{hyp:P}, and let $\mathfrak K_0$ be the barycentric subdivision
of $\mathscr P$. Every balanced network on $\Gamma$ is the boundary network of a relative tropical $2$-cycle on
$\mathfrak K_0$, and the node coefficients of balanced networks form exactly $K$:
\[
\partial Z(\mathfrak K_0)=\mathrm{Bal}(\Gamma),\qquad \rho\bigl(\mathrm{Bal}(\Gamma)\bigr)=K,\qquad\text{hence}\qquad \rho\bigl(Z(\mathfrak K_0)\bigr)=K.
\]
For every triangulation $\mathfrak K$ of $B$ in which $\Gamma$ is full, $\rho(Z(\mathfrak K))\subseteq K$.
\end{mainthm}

\begin{mainthm}[the mirror lift]\label{thm:lift}
With the members $Y_t$ of Theorem~\ref{thm:main}, for every triangulation $\mathfrak K$ of $B'_\sigma$ in which
$\Gamma$ is full and every $z\in Z(\mathfrak K)$ there is an integral $4$-chain $\check\Xi(z)$ on $Y_t$, constructed from
$z$ and finitely many auxiliary choices, with
\[
\partial\,\check\Xi(z)=\sum_q\rho_q(z)\,[\delta_q].
\]
\end{mainthm}

\begin{remark}[comparison with strict tropical cycles]\label{rem:strict-statement}
For a tropical $2$-cycle $(S,j,v)$ in the sense of \cite[Def.~7.2]{CBM}, the chain $\check\Xi(S)$ realises on $Y_t$ the
four-dimensional object $\widetilde S$ that Casta\~no-Bernard and Matessi attach to $S$ in a topological smoothing
\cite[\S7]{CBM} in the following sense. Over the part of $S$ away from the discriminant
it is, up to homotopy in each region, the bundle whose fibre over a point is the two-torus annihilated by $v$ in the dual
torus $\Lambda^\vee\otimes U(1)$. Completing this local construction to a chain with the stated boundary uses
the fibre-torus corrections of Proposition~\ref{lem:F}. Near the discriminant the local pieces agree with those of
\cite{CBM} when the position conditions (C1)--(C2) of Remark~\ref{rem:CBMlift} hold; that remark gives the precise comparison.
\end{remark}

The period computation (Section~\ref{sec:periods}) turns the inclusion $K\subseteq\Rel(\delta)$ given by
Theorems~\ref{thm:real} and~\ref{thm:lift} into equality and gives the nodal locus.

On the resolution side the same tropical objects produce the relations among the exceptional curves, for tropical
$2$-cycles in the sense of \cite{CBM}.

\begin{mainthm}[the resolution side]\label{thm:res}
Let $\Delta$ be admissible with a profile satisfying~\ref{hyp:P}, and let $\widehat X$ be a projective small resolution of
$X'_\sigma$ induced by a regular refinement of $\Sigma'$, and label the node parallelograms by the diagonals of
$\widehat X$. For every tropical $2$-cycle $(S,j,v)$ on $B'_\sigma$ in the sense of \cite[Def.~7.2]{CBM}, with $j$ piecewise
linear, there is an integral $3$-chain $\Xi(S)$ on $\widehat X$, constructed from $(S,j,v)$ and finitely many auxiliary
choices, with
\[
\partial\,\Xi(S)=\sum_q\rho_q(S)\,[C_q].
\]
For any other small resolution the same holds after reversing the sign of $\rho_q$ at every flopped curve. With
Theorem~\ref{thm:lift}, the node coefficients of $S$ are a relation among the exceptional curves of $\widehat X$ and among
the vanishing spheres of $Y_t$ simultaneously.
\end{mainthm}

The construction of the base uses the pair $(\mathcal T,\check h)$, whose existence is a combinatorial question about
the polar polytope.

\begin{mainthm}[unimodular heights]\label{thm:U}
For every admissible reflexive four-polytope $\Delta$, the boundary of $\Delta^\circ$ has a unimodular triangulation
$\mathcal T$ refining its faces, with an integral height $\check h$ such that $\check h$ and $\check h-\check\varphi$ are
strictly convex on the fan over $\mathcal T$. Equivalently, $\Delta^\circ$ has a fine, regular, unimodular triangulation
that is star-shaped from the origin.
\end{mainthm}

The proof is computer-assisted (Appendix~\ref{app:unimodular}). Up to $GL_4(\Z)$ there are $51{,}827$ admissible
reflexive four-polytopes in the classification of Kreuzer and Skarke \cite{KS00}, $40{,}403$ of them without hexagonal
two-faces; for each we exhibit integer heights and verify the conclusion in exact arithmetic. That these are all the
admissible polytopes of the classification rests on the admissibility program of Section~\ref{ssec:utrust}, whose
negative verdicts are only partly replicated by a second implementation.

Finally, the strict tropical $2$-cycles of \cite{CBM} are expected to suffice on their own.

\begin{mainconj}[strict realisation]\label{conj:strict}
Under the hypotheses of Theorem~\ref{thm:real}, the node coefficients of tropical $2$-cycles $(S,j,v)$ on $B'_\sigma$ in the
sense of \cite[Def.~7.2]{CBM}, with $j$ piecewise linear, generate $K$.
\end{mainconj}

By Theorems~\ref{thm:lift} and~\ref{thm:res}, Conjecture~\ref{conj:strict} would realise every relation of $K$ on both
sides by the same objects. Computations find strict tropical $2$-cycles whose node coefficients generate $K$ on the
five examples of Section~\ref{sec:examples} and on $340$ of the $364$ hexagon-free polytopes with $K\neq0$ of the list
of Section~\ref{ssec:list} that satisfy~\ref{hyp:P}. For $X_9$ this is proved in the text (Lemma~\ref{lem:x9belt}), given the triangulation $\mathcal T$ and the triangles of the
cycle recorded in the accompanying data (Appendix~\ref{app:data}). For the other examples the boundary networks and node
coefficients of the cycles, and the shape of the belt cycles (a disc with one field that meets $\Gamma$ only along its
boundary), are checked in the accompanying data (Appendix~\ref{app:data}); for their other cycles the interior conditions of
\cite[Def.~7.2]{CBM} rest on the computations described in Section~\ref{ssec:evidence}, as do the results on the $340$
polytopes. On the other $24$ polytopes the
cycles found generate a sublattice of $K$ of corank between $1$ and $4$, while balanced networks of the shape that such
cycles have along their boundary already generate $K$ (Section~\ref{ssec:evidence}).

\subsection{Scope}\label{ssec:scope}
Hypothesis~\ref{hyp:P} is the only restriction on admissible polytopes in Theorems~\ref{thm:main}--\ref{thm:res}.
It is needed by the method: without it $X'_\sigma$ is not projective over $X$, and the base $B'_\sigma$ is not defined.
It is a genuine restriction. On the list of Section~\ref{ssec:list} ($612$ admissible polytopes up to $GL_4(\Z)$) it holds
for $364$ of the $375$ hexagon-free polytopes with $K\neq0$; of the $142$ polytopes with a hexagonal two-face it holds
for some profile on $83$ ($55$ with every hexagon on the $A_2$ branch, $20$ with every hexagon on the $A_1$ branch, and
$18$ with a mixed profile, with overlaps), and for no profile on $59$. Where it fails there is an explicit obstruction: a nonzero
effective combination of the proper transforms $\tilde f_\epsilon\subset\widehat X$ of curves in the fibres of
$X'_\sigma\to X$ over the del Pezzo points that is homologous in $H_2(\widehat X;\Q)$ to a combination of the
exceptional curves $C_q$ (Proposition~\ref{prop:projective}(3)).

Theorem~\ref{thm:main} is stated for members with real coefficients of the given sign pattern near $\mathcal Z$, for
$t$ small. The classes $[\delta_q]$ are flat sections of the local system $H_3(Y_{t,c};\Z)$ over the members of the
neighbourhood $U$ of item~(3) that lie off the nodal hypersurfaces, and this set is connected; so item~(1) holds for all
of these members, complex coefficients included, with $\delta_q$ at such a member oriented by transport from $c$
(Section~\ref{ssec:relations}).

\emph{Two-faces with one interior lattice point.} Let $G$ be a lattice polygon with primitive edges and exactly one
interior lattice point. Up to $GL_2(\Z)$ and translation, $G$ is the polygon spanned by the primitive ray generators
of the fan of one of the five smooth toric del Pezzo surfaces $S_G$: $\mathbb P^2$, $\mathbb F_1$,
$\mathbb P^1\times\mathbb P^1$, $dP_7$ or $dP_6$. As in Section~\ref{ssec:partialres}, if $G$ is a two-face of $\Delta$
with $\ell(G^\circ)=1$, the point of $X$ on the orbit of the cone over $G$ is the anticanonical cone over $S_G$. For
$\mathbb P^2$ this cone is $\C^3/\mu_3$, with $\mu_3$ acting by scalars, which is rigid \cite{Schlessinger}. For
$\mathbb F_1$ the cone is not smoothable: $G$ has no nontrivial Minkowski decomposition into lattice polygons, and the
base space of the versal deformation is $\operatorname{Spec}\C[\varepsilon]/(\varepsilon^2)$
\cite[(9.1)(iv), (9.2)]{Altmann97}. For $\mathbb P^1\times\mathbb P^1$, $dP_7$ and $dP_6$ the cones are smoothable
\cite[Table~1, Lemma~3.5]{Paper5}; the components of their reduced versal base spaces correspond to the maximal
Minkowski decompositions of $G$ into lattice polygons \cite[(2.5), (8.1)]{Altmann97}. Admissibility therefore keeps
every such type with a smoothable cone point except $\mathbb P^1\times\mathbb P^1$. Its polygon
$\operatorname{conv}\{(\pm1,0),(0,\pm1)\}$ is a lattice parallelogram of twice the area of a unit parallelogram, and
its vertices, placed at height one, span a sublattice of index two; so its cone is the quotient of the conifold by
an involution acting freely off the vertex, not a node. Its only lattice point besides the vertices is the centre,
so it has no subdivision into unit parallelograms and unimodular triangles that contains a parallelogram, and the
nodal partial resolutions used here do not exist for it. The smoothing criterion of \cite[Thm.~4.7]{Paper5} for $X$
itself allows cones over $\mathbb P^1\times\mathbb P^1$; for hexagon-free $\Delta$ satisfying~\ref{hyp:P}, the
relation criterion of Corollary~\ref{cor:smooth} decides the smoothability of $X$ itself (Corollary~\ref{cor:mixed}).
